%%%%%%%%%%%%%%%%%%%%%
% Documento maestro %
%%%%%%%%%%%%%%%%%%%%%
\documentclass{uanl}

%%%%%%%%%%%%%%%%%%%%%%%%%%%%%%%%%%%%%%%%%%%
% Cargando paquetes y definiendo opciones %
%%%%%%%%%%%%%%%%%%%%%%%%%%%%%%%%%%%%%%%%%%%
% Aquí puedes cargar los paquetes que vas a usar. La clase
% fime ya incluye babel, inputenc, graphicx y los de la AMS.
% Cargar un paquete está a tu libertad (y responsabilidad).
\usepackage{hyperref}
    \hypersetup{breaklinks=true,colorlinks=true,linkcolor=black,citecolor=black,urlcolor=black}

%%%%%%%%%%%%%%%%%%%%%
% Definiendo campos %
%%%%%%%%%%%%%%%%%%%%%
\def\titulo{Título de la tesis}
\def\autor{Juan Jesús Torres Solano}
\def\matricula{0000000}
\def\grado{Maestría en Ciencia de Datos}
% En caso de que el grado tenga orientación o especialidad llenar el siguiente
% campo dejando un ESPACIO INICIAL, en caso contrario, dejar vacío
%\def\orientacion{Con Orientación en Fibras Ópticas, Fotónica y Sensores Ópticos.}
% Coloca el mes con mayúscula inicial
\def\fecha{marzo 2025}

\def\asesor{Dr. José Anastasio Hernández Saldaña}
\def\revisorA{Nombre del revisor A}
\def\revisorB{Nombre del revisor B}
% En el caso de que tu tesis sea de doctorado activa la variable cambiándola a \doctoradotrue
% y define tus otros dos revisores
\newif\ifdoctorado\doctoradotrue
\def\revisorC{Nombre del revisor C}
\def\revisorD{Nombre del revisor D}
% El visto bueno siempre va
\def\vobo{Dra. Azucena Yoloxóchitl Ríos Mercado}

%%%%%%%%%%%%%%%%%%%%%%%
% Inicia el documento %
%%%%%%%%%%%%%%%%%%%%%%%
\begin{document}

\frontmatter
\pagestyle{main}

%%%%%%%%%%%%%%%%%%%%%%%%
% Primer portada: UANL %
%%%%%%%%%%%%%%%%%%%%%%%%
\thispagestyle{empty}

\begin{scshape}
\begin{center}
	{\Large \uanl} \\[5mm]
	{\large \fcfm} \\[5mm]
	{\large \mcd}
	\vskip15mm
	\includegraphics[height=55mm]{uanl}
	\vskip12mm
	\begin{tabular}{p{11cm}}
		\centering
		{\large \titulo}
	\end{tabular}
	\vskip7mm
	{por}\\[7mm]
	{\large \autor}\\[7mm]
        % Si en tu posgrado la tesis no es opcional (como sí lo es en licenciatura), no modifiques esta línea:
	{como requisito parcial para obtener el grado de}\\[3mm]
	\MakeUppercase{\grado}\\
	\orientacion
	\vfill
	\fecha\\
	\uanl\\
	\fcfm
\end{center}
\end{scshape}

%%%%%%%%%%%%%%%%%%%%%%%%%%%%%
% Carta del comité de tesis %
%%%%%%%%%%%%%%%%%%%%%%%%%%%%%
\newpage
\thispagestyle{empty}
\enlargethispage{5mm}

{\renewcommand{\baselinestretch}{1.1}\selectfont
\begin{center}\vspace*{-25mm}\hspace*{-10mm}
\begin{minipage}{170.5mm}
\hspace{-1.5mm}\includegraphics[height=20mm]{uanl}\hfill\raise0mm\hbox{\includegraphics[height=20mm]{fcfm}}
\hrule\vspace{0.5mm}
\scalebox{.5}{\MakeUppercase{\uanl}}\hfill\scalebox{.5}{\MakeUppercase{\fcfm}}\medskip
\end{minipage}
%\vskip4mm{\sc\large\uanl\\\fcfm\\[3pt]\pifi}\vskip6mm
\end{center}

Los miembros del Comité de Tesis recomendamos que la Tesis <<\titulo>>, realizada por el alumno \autor, con número de matrícula \matricula, sea aceptada para su defensa como requisito parcial para obtener el grado de \grado\orientacion.
\ifdoctorado\vskip10mm\else\vskip8mm\fi

\begin{center}
El Comité de Tesis\\
\ifdoctorado\vskip15mm\else\vskip25mm\fi

\ifdoctorado{%%%
\begin{tabular}{p{37mm}p{21mm}p{12mm}p{21mm}p{37mm}}
	\cline{2-4}
	& \multicolumn{3}{c}{\asesor} & \\
	& \multicolumn{3}{c}{Asesor}  & \\[15mm]
	\cline{1-2} \cline{4-5}
	%\multicolumn{2}{c}{\revisorA} & & \multicolumn{2}{c}{\revisorB} \\
	%\multicolumn{2}{c}{Revisor}   & & \multicolumn{2}{c}{Revisor}   \\[17mm]
	\cline{1-2} \cline{4-5}
	\multicolumn{2}{c}{\revisorC} & & \multicolumn{2}{c}{\revisorD} \\
	\multicolumn{2}{c}{Revisor}   & & \multicolumn{2}{c}{Revisor}   \\[2mm]
	& \multicolumn{3}{c}{Vo. Bo.} & \\[14mm]
	\cline{2-4}
	& \multicolumn{3}{c}{\vobo}   & \\
	& \multicolumn{3}{c}{Subdirector de Estudios de Posgrado}   & \\ &&&&
\end{tabular}
}\else{%%%
\begin{tabular}{p{37mm}p{21mm}p{12mm}p{21mm}p{37mm}}
	\cline{2-4}
	& \multicolumn{3}{c}{\asesor} & \\
	& \multicolumn{3}{c}{Asesor}  & \\[19mm]
	\cline{1-2} \cline{4-5}
	\multicolumn{2}{c}{\revisorA} & & \multicolumn{2}{c}{\revisorB} \\
	\multicolumn{2}{c}{Revisor}   & & \multicolumn{2}{c}{Revisor}   \\[2mm]
	& \multicolumn{3}{c}{Vo. Bo.} & \\[17mm]
	\cline{2-4}
	& \multicolumn{3}{c}{\vobo}   & \\
	& \multicolumn{3}{c}{Coordinadora de la Maestría en Ciencia de Datos}   & \\ &&&&
\end{tabular}
}\fi%%%

\vfill

\snnl, \MakeLowercase{\fecha}

\end{center}}

% Dedicatoria

\thispagestyle{empty}
\vspace*{17mm}

\begin{flushright}
\begin{itshape}

Aquí puedes poner tu dedicatoria\\
si es que tienes una.\bigskip\bigskip

\end{itshape}
\end{flushright}



\tableofcontents
\listoffigures
\listoftables

%Agradecimientos

\chapter{Agradecimientos}
\markboth{Agradecimientos}{}

Aquí puedes poner tus agradecimientos. (No olvides agradecer a tu comité de tesis, a tus profesores, a la facultad y a CONACyT en caso de que hallas sido beneficiado con una beca).

%Resumen

\chapter{Resumen}
\markboth{Resumen}{}

{\renewcommand{\baselinestretch}{1.1}\selectfont
{\setlength{\leftskip}{10mm}
\setlength{\parindent}{-10mm}

\autor.

Candidato para obtener el grado de \grado\orientacion.

\uanl.

\fcfm.

Título del estudio: \textsc{\titulo}.

\noindent Número de páginas: \pageref*{lastpage}.}

%%% Comienza a llenar aquí
\paragraph{Objetivos y método de estudio:}
Aquí debes poner tus objetivos y métodos de estudio. (Este es el formato).

\paragraph{Contribuciones y conlusiones:}
Y aquí tus contribuciones y conclusiones. (También es parte del formato).

\bigskip\noindent\begin{tabular}{lc}
\vspace*{-2mm}\hspace*{-2mm}Firma del asesor: & \\
\cline{2-2} & \hspace*{1em}\asesor\hspace*{1em}
\end{tabular}}



\mainmatter
\pagestyle{uanl}

%%% Haz un documento para cada capítulo
\chapter{INTRODUCCIÓN}
Durante las etapas iniciles de proyectos de ingenieria civil, vial, petrolera, geohidrologica, se requiere conocer el medio sobre el cual se desarrollaran los proyectos, ya sea para establecer medidad de accion de excabacion, aprobechamiento o explotacion de algun mineral, para ello se requiere trabajar 
que es la geofisica:

La exploración geofísica consiste en un conjunto de metodologías que a través de la medición de propiedades petrofisicas del subsuelo es 

eL geoelectrica

para que sirve


oportunidades de la cuiencia de datos en el análisis en tiempo real de la respuesta 




\chapter{DELIMITACIÓN Y PLANTEAMIENTO DEL PROBLEMA}

Durante un trabajo de prospección geofísica, al realizar adquisición de datos geoeléctricos \textit{in situ}, no es posible conocer el resultado del trabajo hasta una vez realizado procesos de inversión de datos resultado de prospección, por lo que no se tiene certeza de si una lectura presenta inconsistencias o si forma parte de una respuesta esperada para determinado resultado.\\

La problemática radica en que el muestreo planeado, en la etapa de planeación de adquisición, puede no ser efectivo o no reflejar una distribución esperada, debido a que el medio que se prospecta carece de homogeneidad, siendo solo evidente durante la exploración directa del medio, lo que en muchos casos no se logra a identificar, sumando ambigüedad al proceso de interpretación.\\ 

Como primera aproximación a una solución se considerara un primer supuesto, en el cual se cuenta con respuestas de sondeos eléctricos verticales, más simple pero no menos complejo de interpretar, al contar con un solo segmento de señal, dentro de esta señal podemos encontrar distintas unidades geológicas, profundidad de acuífero, espesores de unidades y profundidad de exploración.\\

La propuesta de solución es mediante Aprendizaje automático (ML, por sus siglas en inglés), empleando una técnica de aprendizaje que permita identificar patrones que estén asociados a la respuesta de un modelo de inversión, de esta manera obtener un modelo pronóstico de la inversión, sirviendo de guía para incrementar la densidad de muestreo para mejorar el modelo y respuestas.

\chapter{JUSTIFICACIÓN}
Existen múltiples aplicaciones de ML y DL en el procesamiento, modelado e interpretación geofísica \citep{li2024, liu2020,el2001,wrona2018}, algunas de estas aplicaciones corresponden a implementaciones académicas, al igual es posible identificar aplicaciones comerciales implementadas en software principalmente de exploración sísmica de hidrocarburos \citep{diaferia2024high, panebianco2024automated}.\\


%Es de esperar que al implementar ciencias de datos en diversas áreas de la ciencia, se optimicen aspectos básicos, mejorando la analítica e interpretación de resultados, o bien, modelando predicciones de tendencias.

La aplicación de herramientas de ML durante la ejecutar de muestreo de Sondeo Eléctrico Vertical (SEV), en particular durante el proceso de adquisición \textit{in situ}, presenta la posibilidad de evaluar y ampliar el intervalo de muestreo original, mejorando la adquisición tradicional, permitiendo evaluar el muestreo contra la regresiones y establecer nuevos intervalos de oportunidad.\\

Durante el muestreo tradicional los intervalos se diseñan considerando el objetivo de exploración (idealmente), este análisis previo es un factor determinante, en esta etapa se establece el intervalo de muestreo que permitiría identificar el objeto de exploración, este objeto de exploración pueden ser unidades geológicas, acuíferos, fallas, zonas de fracturas, estructuras antropogénicas, etc.\\

De manera general se busca mantener un intervalo de muestreo menor a la frecuencia de ocurrencia del objetivo de estudio, por lo que el éxito de la exploración dependerá en su totalidad de la planeación previa de la adquisición, lo que implicaría conocer previamente la conformación, distribución y espesor de cada unidad, lo que resulta irrisorio.\\

De manera que un modelo entrenado permita generar múltiples regresiones e identificar regiones de interés no cubiertas por el muestreos y de esta manera generar puntos adicionales \textit{in situ} optimizando la adquisición de datos e impactando en la calidad del muestreo , mejorando el acotamiento de las frecuencias identificadas.\\

Pese a existir aplicaciones de ML implementadas en geofísica, no se identifica alguna enfocada este esta problemática en particular, sin embargo hay aplicaciones en otros campos de estudio con un enfoque en el muestreo y clasificación de datos no paramétricos \citep{entezami2022non, bkassiny2013multidimensional, shi2021non} empleando técnicas como Dirichlet Process Mixture Model (DPMM), radial basis function network (RBFN),  Multiple Point Statistics (MPS) y Bayesian Compressive Sensing (BCS).\\


porque random forest y su copmaracion con BCS

\chapter{Conclusión}

\section{Estado actual}

\section{Perspectiva a futuro}

\appendix
%%% Haz un documento para cada apéndice, si es que tienes
\chapter{Apéndice I}






\backmatter
\pagestyle{main}

% Aquí va la bibliografía, puedes usar el entorno de LaTeX (thebibliography)
% o la herramienta BibTeX. En caso de que optes por BibTeX, puedes usar
% alguno de los archivos de estilo (mighelbib.bst o mighelnat.bst) incluidos
% en el paquete, cuyos diseños armonizan con el diseño de tesis provisto por
% fime.cls. Para muestra, basta un botón:
\bibliographystyle{mighelnat}
\bibliography{MiBiblio}

\label{lastpage}
%Autobiografia

\chapter*{Resumen autobiográfico}
\thispagestyle{empty}

\begin{center}
\autor

Candidato para obtener el grado de\\
\grado\\
\orientacion\bigskip

\uanl\\
\fcfm\bigskip

Tesis:\\
\textsc{\large\titulo}
\end{center}\bigskip

%Aquí va tu historia
Aquí va tu historia. Recuerda que debe incluir: lugar y fecha de nacimiento, nombre de los padres, escuelas y universidades en las que se graduó después de la preparatoria, títulos o grados obtenidos (no incluir los estudios que se están concluyendo), experiencia profesional y organizaciones profesionales a las que pertenece (no incluir lista de publicaciones).


\end{document}

